\section{The once-punctured CM curve: genuine actions, graded images, and inertia}
\label{sec:G}

\subsection{The filtration and the first correction to the question}

Throughout this part we put $K=\mathbb Q(i)$, $E:Y^2=X^3-X$,
$U=E-\{O\}$, $P=\pi_1(U_{\overline K},\vec v)^{(5)}$,
$N_r=P/\gamma_{r+1}P$, and $\Delta=G_{K^{\rm ab}}$.
The base point $\vec v$ is the rational tangent defined by $t=X/Y$.
We reserve $x,y$ for a symplectic CM eigenbasis in homology, or for its
formal Lie lifts, and write $\omega=[x,y]$ for the Lie boundary. The actual
inertia generator is denoted $c$; a symplectic expansion sends
$\log c$ to $\omega$. These conventions avoid identifying the group
commutator with the bracket of logarithms. The two CM characters are
$\psi,\overline\psi:G_K\longrightarrow\mathbb Z_5^\times$;
$\psi\overline\psi=\chi_5$.

\begin{statement}{G.1}{Refuted}{An arbitrary homogeneous term is not an intrinsic Johnson homomorphism}
The expression ``the degree-$m$ part of the action of every
$\delta\in G_1$'' does not, without an expansion, define a map to the
degree-$m$ special derivations. The intrinsic comparison is
\[
 j_m:A_m/A_{m+1}\longrightarrow
 \operatorname{Hom}_{\mathbb Z_5}(V,L_{m+1}),\qquad
 j_m(\alpha)(\bar u)=\alpha(u)u^{-1}\pmod{\gamma_{m+2}P},
\]
where $A_m=\ker(\operatorname{Aut}(P)\to\operatorname{Aut}(N_m))$,
$V=P^{\rm ab}$ and $L_r=\gamma_rP/\gamma_{r+1}P$.
It is a continuous injective homomorphism of compact additive
$\mathbb Z_5$-modules. On the subgroup fixing $c$ its image consists of
derivations annihilating $\omega$.
\end{statement}
\begin{proof}
Changing a lift $u$ by an element of $\gamma_2P$ changes the displayed
error only in degree at least $m+2$. The commutator identities show that
the resulting map is additive in $\bar u$, obeys the derivation rule on
iterated brackets, and is additive in $\alpha$ modulo $A_{m+1}$.
Its vanishing is exactly the condition $\alpha\in A_{m+1}$. Applying
$\alpha$ to the boundary gives $j_m(\alpha)(\omega)=0$.
For an element outside $A_m$, the displayed error is not even an element
of $\gamma_{m+1}P/\gamma_{m+2}P$. Rationally one may instead choose a
filtered isomorphism of the Malcev Lie algebra with
$\widehat{\mathbb L}(V\otimes\mathbb Q_5)$ and take homogeneous
components of $\log\alpha$. Those components change under changes of
expansion, and the Campbell--Hausdorff formula introduces brackets of
lower components in their composition. Thus the unrestricted notation
in the question requires this additional choice. Below, graded image
always means $j_m(G_K\cap A_m)$; homogeneous logarithms are used only
after the expansion has been specified.
\end{proof}

\begin{statement}{G.2}{Proved}{CM symmetry and equality of the two rational graded images}
There is a genuine order-four automorphism $J$ of $P$, commuting with
the tangential $G_K$-action, induced by $[i]$ on $E$.
There exists a $J$-equivariant rational symplectic expansion. In such an
expansion all odd homogeneous components of $\log\rho_s(G_1)$ vanish,
and the degree-$m$ component has only characters
$\psi^a\overline\psi^b$ with $a+b=m$ and $a\equiv b\pmod4$.
Intrinsically, for every $m>0$,
\[
 \operatorname{span}_{\mathbb Q_5}j_m(\Delta\cap\rho_s^{-1}(A_m))
 =\operatorname{span}_{\mathbb Q_5}j_m(G_1\cap\rho_s^{-1}(A_m)).
 \tag{G.1}
\]
The equality is of rational subspaces; it asserts no equality of integral
lattices.
\end{statement}
\begin{proof}
The local parameter satisfies $[i]^*t=it$. On the covers
$t_n^{5^n}=t$, use $t_n\mapsto it_n$. These maps commute with
$t_{n+1}^5=t_n$, have order four, and are defined over $K$.
They give a $G_K$-invariant path between the corresponding tangential
fibres, hence the asserted genuine lift. It fixes local inertia.
In the rational Malcev Lie algebra, average a linear lift of $V$ over
$\langle J\rangle$ to obtain equivariant free generators. To straighten
the boundary, successively remove its homogeneous errors. In degree
$n+2$ the change of generators is governed by
$(u,v)\mapsto[u,y]+[x,v]$ from $L_{n+1}^2$ to $L_{n+2}$.
This map is surjective: the Jacobi identity writes every bracket of
length $n+2$ as a sum of brackets of a generator with a word of length
$n+1$. Average each correction over $J$; its target error is $J$-fixed.
The corrections converge in the degree filtration and give the required
expansion. Conjugation by $J$ acts on a character $(a,b)$ by $i^{a-b}$.
Commutation therefore forces $a-b\equiv0\pmod4$, and in particular
$m=a+b$ is even. The same argument with $J^2$ gives integral vanishing
of every odd Johnson image because $2$ is a unit in $\mathbb Z_5$.

For (G.1), the quotient $G_K/G_1$ is abelian by CM, so for
$g\in G_1\cap\rho_s^{-1}(A_m)$ and $t_0\in G_K$ the element
$[t_0,g]$ belongs to $\Delta\cap\rho_s^{-1}(A_m)$, and
\[
 j_m([t_0,g])=(T_{t_0}-1)j_m(g).
\]
Choose $t_0$ to be a Frobenius at a good prime away from $5$.
Every eigenvalue of $T_{t_0}$ on this degree is a product of $m$
homological Frobenius eigenvalues, and has complex absolute value
$q^{m/2}>1$. Thus $T_{t_0}-1$ is invertible over $\mathbb Q_5$.
The span on the right of (G.1) is $T_{t_0}$-stable, so the displayed
commutators span it. The reverse inclusion follows from $\Delta\subset
G_1$. This is the required comparison, rather than a map assigning
independent Johnson values to all elements of $G_1$ in every degree.
\end{proof}

\subsection{Integral special derivations and the degree-two Kummer character}

\begin{statement}{G.3}{Proved}{The exact kernel presentation of special derivations}
For $R=\mathbb Z_5$ or $\mathbb Q_5$, let $\mathfrak s_m(R)$ be the
degree-$m$ derivations of $\mathbb L_R(x,y)$ that annihilate $\omega$.
There is an exact sequence of finite free $R$-modules
\[
 0\longrightarrow\mathfrak s_m(R)\longrightarrow L_{m+1}^2
 \xrightarrow{\ (u,v)\mapsto[u,y]+[x,v]\ }L_{m+2}
 \longrightarrow0.                                      \tag{G.2}
\]
In particular $\mathfrak s_2(\mathbb Z_5)=\mathbb Z_5\operatorname{ad}\omega$.
The latter derivation is central in the whole special derivation algebra.
\end{statement}
\begin{proof}
A derivation is uniquely determined by $D(x)=u,D(y)=v$. The boundary
equation gives the kernel in (G.2); surjectivity was proved in G.2 and is
integral. The target is free, so the sequence splits as modules.
In degree three the basis is $[x,\omega],[y,\omega]$. Substitute arbitrary
linear combinations for $D(x),D(y)$. Independence of the three degree-four
Lyndon brackets gives
$D(x)=a[\omega,x]$ and $D(y)=a[\omega,y]$.
Finally $[D,\operatorname{ad}\omega]=\operatorname{ad}(D\omega)=0$.
\end{proof}

\begin{statement}{G.4}{Known}{The discriminant coefficient for the specified tangent}
For the tangent $t=X/Y$ and $\delta\in G_1$, the genuine action has
degree-two term
\[
 j_2(\rho_s(\delta))=a_s(\delta)\operatorname{ad}\omega,
 \qquad
 a_s=\frac1{12}\kappa_{64}=\frac12\kappa_2=\kappa_{1+i}
 \quad\hbox{on }G_1.                                    \tag{G.3}
\]
Here $\kappa_a:G_K\to\mathbb Z_5(1)$ is the continuous Kummer
cocycle specified by compatible $5$-power roots of $a$; its restriction
to $G_1$ is independent of those roots and is a homomorphism.
Moreover $a_s(\Delta)=\mathbb Z_5$.
\end{statement}
\begin{proof}
We apply \citet[Theorem 6.10.3]{Nakamura2013Eisenstein}, not merely its
abstract. In his normalization $z=[x,y]^{-1}$ and the constant term of
the metabelian invariant $E_\sigma$ is $\kappa_{\mathcal D}(\sigma)/12$.
His equation (6.10.9) gives $E_{\operatorname{Int}z}=-1$; hence
$E_{\operatorname{Int}c}=1$. By G.3 this constant term is exactly the
coefficient of $\operatorname{ad}\omega$. Put $y_{\rm Nak}=2Y$ in his
equation $y_{\rm Nak}^2=4X^3-g_2X-g_3$. Then $g_2=4,g_3=0$,
$\mathcal D=64$, and his tangent is $-X/Y$. The ratio $-1$ has compatible
$5$-power roots in $K$, so it changes no Kummer class. This proves the
first equality. Finally $(1+i)^2=2i$ and $i$ has compatible $5$-power
roots in $K$, which proves the remaining equalities, with integral
$5$-adic coefficients.

For the final assertion put $F=K(\mu_{5^\infty})$ and
$M=F((1+i)^{1/5^\infty})$. At the prime $(1+i)$, the extension $F/K$
is unramified and $v(1+i)=1$, so $\operatorname{Gal}(M/F)=\mathbb Z_5$.
Conjugation on this group by $G_K$ is multiplication by $\chi_5$.
The cyclotomic image contains a value whose difference from $1$ is a
unit. Its commutators with $\operatorname{Gal}(M/F)$ therefore generate
all of $\mathbb Z_5$. The image of the closed commutator subgroup
$\Delta$ is consequently all of that group. This proves the assertion
on the actual absolute Galois subgroup, without replacing it by the
cyclotomic kernel.
\end{proof}

\begin{statement}{G.5}{Proved}{Every section in the cusp packet}
Orient packet cocycles by $c$ and write
$s_b(\sigma)=c^{b(\sigma)}s(\sigma)$, with
$b\in Z^1(G_K,\mathbb Z_5(1))$. Then
\[
 a_{s_b}=\kappa_{1+i}+b\quad\hbox{on }G_1.                \tag{G.4}
\]
There is precisely one packet class with $a_{s_b}|_\Delta=0$:
$[b]=-[\kappa_{1+i}]$. It is represented by a rational tangent change.
Every other pro-$5$ packet class has nonzero degree-two character on
$\Delta$.
\end{statement}
\begin{proof}
Conjugation by $c^{b(\sigma)}$ adds $b(\sigma)\operatorname{ad}\omega$
in degree two, proving (G.4). If a cocycle with values in
$\mathbb Z_5(1)$ restricts trivially to $\Delta$, it descends to the
abelian group $G_K^{\rm ab}$. For commuting $g,h$ its cocycle values
satisfy
$(\chi_5(g)-1)b(h)=(\chi_5(h)-1)b(g)$.
Choose $g$ with $\chi_5(g)-1$ a unit. The equality says exactly that
$b$ is the coboundary of $b(g)/(\chi_5(g)-1)$. Thus restriction of
$H^1(G_K,\mathbb Z_5(1))$ to $\Delta$ is injective. Apply this to
$[\kappa_{1+i}]+[b]$. The inverse Kummer class is represented by
$(1+i)^{-1}$, so the exceptional class lies in the rational tangent
packet as well. The coordinate sign in the tangent rescaling depends on
whether local positive inertia is called $c$ or $c^{-1}$; formula
(G.4) fixes the convention completely.
\end{proof}

\subsection{Degree four: two vanishing components and one nonzero component}

Let $b_w$ denote the standard Lyndon bracketing of a word $w$ in $x<y$.
Put $F(x)=y,F(y)=0$, and define
\[
 \begin{aligned}
 \epsilon_4(x)&=\operatorname{ad}_x^4y,\\
 \epsilon_4(y)&=[y,\operatorname{ad}_x^3y]
                -[[x,y],\operatorname{ad}_x^2y],\qquad
 T=[F,\epsilon_4].
 \end{aligned}                                             \tag{G.5}
\]

\begin{statement}{G.6}{Refuted}{The three claimed nonzero degree-four components}
The special degree-four module has rank three and characters
$(3,1),(2,2),(1,3)$, as suggested in the prompt. For this CM curve,
however, the first and third components of the Galois image vanish:
\[
 b_+=b_-=0,
 \qquad D_4(\delta)=b_0(\delta)T.
 \tag{G.6}
\]
Here $T$ is primitive over $\mathbb Z_5$, with
\[
 T(x)=-2b_{xxxyy}-b_{xxyxy},\qquad
 T(y)=2b_{xxyyy}+3b_{xyxyy}.                              \tag{G.7}
\]
In particular the assertion that each of $b_+,b_0,b_-$ is nonzero on
$\Delta$ is false.
\end{statement}
\begin{proof}
The multigraded kernel of (G.2) has rank one at each of the three
displayed characters and zero elsewhere. Substitution in (G.5) gives
(G.7), $T(\omega)=0$, and the weights of $\epsilon_4,T,F^2\epsilon_4$
are respectively $(3,1),(2,2),(1,3)$. The last vector is divided by its
factor $2$ to obtain a primitive basis. These divisions are units at
$5$. The action of $J$ on the outer two lines is $-1$, whereas it is
$1$ on the middle line. G.2 therefore kills the outer two arithmetic
components. This refutation follows from an actual geometric CM
automorphism, and does not depend on a dimension heuristic.
The program \texttt{g\_special.py} verifies the boundary equations,
all Lyndon coordinates in (G.7), and the ranks over $\mathbb Q$.
The normalization (G.5) agrees, up to the stated sign convention,
with \citet[(3.1.1)]{Tsunogai1995}. The kernel calculation here
uses the full special module, independently of relations in the
smaller algebra generated by these derivations.
\end{proof}

\begin{statement}{G.7}{Known}{The surviving degree-four character and its field}
In the normalization (G.5), let $\kappa_{11}$ be Nakamura's character
defined using the same symplectic eigenbasis. Then on $G_1$,
\[
 b_0=\frac{\kappa_{11}}{48}.                             \tag{G.8}
\]
To compare with the lattice-derived eigenbasis in
\citet{IshiiGoto2026}, let $d\in\mathbb Z_5^\times$ be its
determinant relative to our chosen symplectic basis, and normalize
its second eigenvector by $d^{-1}$. Then
$12\kappa_{11}=d\chi^{\rm ell}_{(2,2)}=-12d\chi^{(4)}_2$,
where $\chi^{(4)}_2$ is their conductor-four Soul\'e cocycle
(Definition 2.14). For the projector basis of $(\Omega,i\Omega)$,
the change-of-basis matrix has columns
$(1,-i_5)/2,(1,i_5)/2$, so $d=i_5/2$ in that orientation.
The restriction of $b_0$ to $\Delta$
is nonzero. Formula (G.8) specifies a class in
$H^1(K,\mathbb Q_5(2))$ (with the usual integral lattice after the
displayed unit factors), not three independent CM classes.
\end{statement}
\begin{proof}
We spell out the comparison with Nakamura's invariant. In the
metabelian Lie algebra the commutator ideal is
$\mathbb Q_5[X_0,Y_0]\omega$, where $X_0=\operatorname{ad}x$ and
$Y_0=\operatorname{ad}y$. A special derivation has images
$X_0h\omega,Y_0h\omega$, and Nakamura's $\alpha$-coordinate is $-h$
because his boundary generator is $\omega^{-1}$ at group level.
For $T$, (G.7) gives $h=2X_0Y_0$. In
\citet[Theorem A, equivalently Corollary 4.12]{Nakamura1995}, the
coefficient of $X_0Y_0$ is
$\kappa_{11}/(1-5^2)$. Thus $-2b_0=-\kappa_{11}/24$.
Under a change of eigenbasis by factors $u,v$, reindexing the
theta-product defining $\kappa_{11}$ multiplies its coefficient
by $(uv)^{-1}$. The normalization $u=1,v=d^{-1}$ therefore
gives the factor $d$. One must also apply the published
\citet[p.~492, Errata to N95, correction of (3.11.4)]{Nakamura2013Eisenstein}:
the exponent in the Kummer property is $12\kappa_{ij}$.
Thus the raw theta-product character of \citet[Definition 1.3]{IshiiGoto2026}
is twelve times Nakamura's coefficient character, before the
basis correction. Remark 1.4 of that paper identifies the basis
comparison but does not display this erratum factor.
Their Theorem 1.7, at
$w_K=4,m=2,B_2=1/6,p=5$, gives
$\chi^{\rm ell}_{(2,2)}=-12\chi^{(4)}_2$, proving
$\kappa_{11}=-d\chi^{(4)}_2$ and hence
$b_0=-d\chi^{(4)}_2/48$. Keeping $d$ and the factor $12$
explicit prevents a normalization
choice from becoming an asserted canonical equality.
Nonvanishing is the published
\citet[Corollary 4.10]{Ishii2025}; its assumptions are class number one,
CM by $\mathbb Z[i]$, and the split prime $5$, all satisfied here.
The character occurs in degree four because the metabelian comparison
in degree four is an isomorphism on the three special lines: their
images are the three independent monomials of degree two. Nonvanishing
on $\Delta$ follows from G.2, or directly by applying a commutator with
a Frobenius whose $\chi_5^2$-value is not $1$.
\end{proof}

The requested polylogarithmic comparison was also checked in the
primary source: \citet[Theorems 4.19 and 5.7, the latter also
Theorem 0.4]{BannaiKobayashiTsuji2010} describe the syntomic elliptic
polylogarithm and its specialization at torsion points of order
prime to $p$ by $p$-adic Eisenstein--Kronecker numbers. Their
hypotheses are $p\ge5$, unramified in the CM field, and good
reduction; they allow supersingular reduction. That specialization
formula alone asserts no global nonvanishing for the present
Galois class. Formula (G.8) instead uses the explicit étale
Kummer comparison and its published nonvanishing theorem, with
all basis factors displayed.

\subsection{Fixed subgroups for every pro-5 cusp section}

\begin{statement}{G.8}{Proved}{Constants of a special derivation detected metabelianly}
Let $k$ be a characteristic-zero field, and let $D$ be a positive-degree
multihomogeneous special derivation of $\mathbb L_k(x,y)$.
If its induced derivation on the metabelian quotient is nonzero, then
\[
                         \ker D=k\omega.                \tag{G.10}
\]
In particular this holds for the middle derivation $T$ of (G.5) and
for each even $m\ge4$ for Tsunogai's derivation
\[
 \epsilon_m(x)=\operatorname{ad}_x^m y,\qquad
 \epsilon_m(y)=\sum_{r=0}^{m/2-1}(-1)^r
    [\operatorname{ad}_x^r y,\operatorname{ad}_x^{m-1-r}y].
 \tag{G.9}
\]
Thus even one such derivation suffices for the fixed-subgroup
argument; an unbounded family will be used separately for normalizers.
\end{statement}
\begin{proof}
Put $V_j=\operatorname{ad}_x^j y$. The ideal generated by $y$ is
freely generated as a Lie algebra by $V_0,V_1,\ldots$: the semidirect
product of that free Lie algebra with $kx$, with $[x,V_j]=V_{j+1}$,
has the universal property of $\mathbb L_k(x,y)$.
The metabelian commutator ideal is $k[X_0,Y_0]\omega$ with
$X_0=\operatorname{ad}x$, $Y_0=\operatorname{ad}y$. If $D$ has
character $(a,b)$, the special equation gives
\[
 D(x)=X_0h\omega,\quad D(y)=Y_0h\omega
 \quad\text{metabelianly},\qquad
 h=\lambda X_0^{a-1}Y_0^{b-1},\quad \lambda\ne0.
\]
Indeed $X_0D(y)=Y_0D(x)$ and $X_0,Y_0$ are relatively prime.
The nonzero special character has $a,b\ge1$: a pure-letter Lie
word of length greater than one vanishes, and a boundary character
with $a=0$ or $b=0$ has zero special kernel. If $a=b=1$, G.3
makes $D$ a nonzero multiple of $\operatorname{ad}\omega$; the
centralizer argument in G.10 proves (G.10). In every other case,
exchange $x,y$ if necessary so that $a\ge2$.

A Lie monomial in the $V_j$ contributes nontrivially to the
metabelian quotient only if it involves at most one $V_j$ with
$j\ge1$. Consequently the nonzero coefficient $\lambda$ gives
exactly the following largest-index terms:
\[
\begin{aligned}
 D(x)&=\lambda\operatorname{ad}_{V_0}^{b-1}V_{a+1}
       +\text{terms with indices at most }a,\\
 D(V_0)&=\lambda\operatorname{ad}_{V_0}^{b}V_a
       +\text{terms with indices at most }a-1.
\end{aligned}
\]
There is only one possible Lie monomial with the indicated largest
letter and the remaining letters all $V_0$, so no coordinate
ambiguity is involved. Since $D(V_1)=D(\omega)=0$, induction with
$D(V_{j+1})=[D(x),V_j]+[x,D(V_j)]$ yields, for $j\ge2$,
\[
 D(V_j)=\lambda[\operatorname{ad}_{V_0}^{b-1}V_{a+j-1},V_1]
       +\text{terms with indices at most }a+j-2.          \tag{G.11}
\]

We record explicitly the no-cancellation argument. For a polynomial
in the free associative algebra on the $V_j$, replacing each
occurrence of $V_j$ in turn by a new letter $Z$ gives a nonzero
formal derivative whenever the polynomial involves $V_j$: a word
with one marked occurrence reconstructs its original word uniquely.
Substitution $Z\mapsto[W,V_q]$ is injective on polynomials linear
in $Z$, provided $W$ is new. Write such a polynomial as $\sum uZv$,
combining identical pairs, and choose the least length of a left
word $u$ with nonzero coefficient. In
$\sum u(WV_q-V_qW)v$, the occurrences of $W$ with that least
left length come only from the first terms and cannot cancel.
The same proof works for $[V_q,W]$. Compositions of these
substitutions remain injective.

Let $w$ be a nonzero homogeneous constant of degree greater than
one, and suppose its expression in the $V_j$ has largest index
$j\ge2$. The new letter $V_{a+j-1}$ has index greater than $j$.
In $D(w)$ it arises only by the substitution in (G.11) for $V_j$.
This is a composition of the injective substitutions just proved,
so its contribution is nonzero. Thus a constant belongs to
$\mathbb L(V_0,V_1)$. The largest-index term of $D(V_0)$ then
excludes $V_0$ by the same argument. Hence $w\in kV_1=k\omega$.
In degree one, the $V_{a+1}$ term of $D(x)$ cannot cancel with
$D(y)$, and $D(y)\ne0$, so no nonzero linear constant exists.
Every nonhomogeneous constant decomposes into homogeneous constants,
which proves (G.10).

For $\epsilon_m$, the metabelian polynomial is $h=X_0^{m-2}$;
for $T$ it is $h=2X_0Y_0$, as checked in G.7. Direct telescoping
in (G.9) verifies the special equation. In the notation above,
$T(x)=2[V_0,V_3]+3[V_1,V_2]$ and
$T(V_0)=2[V_0,[V_0,V_2]]-[V_1,[V_0,V_1]]$, in agreement with
(G.11). These give both advertised special cases.
\end{proof}

\begin{statement}{G.9}{Known}{Unbounded extreme arithmetic derivations}
For every cusp-packet section, the rational graded image of $\Delta$
contains the line $\mathbb Q_5\epsilon_{4k+2}$ for every $k\ge1$.
The same assertion holds at a split prime $\ell\ge5$, replacing these
degrees by $m=2+k(\ell-1)$.
\end{statement}
\begin{proof}
By \citet[Theorem 1.5(4)]{Ishii2025}, the character of index
$(1+k(\ell-1),1)$ is nonzero (in fact surjective). The index satisfies
the CM condition because $4\mid\ell-1$ when $\ell$ splits in
$\mathbb Q(i)$. Nakamura's metabelian comparison identifies it with
the extremal character $(m-1,1)$ in degree $m$.
Its restriction to the $m$th outer filtration has nonzero rational
image: otherwise it would factor through a finite succession of
strictly smaller weights, and equivariance under a good Frobenius
would give a nonzero map between unequal weights. This is also the
argument in \citet[Claim 4.14.1 and equation 4.14.2]{Nakamura1995}.
The ambient special space of that extremal character has dimension
one by (G.2), and (G.9) spans it. The CM torus separates its character
from all other characters. G.2 transfers its rational span to
$\Delta$.

One must still remove the inner degree-two term of a genuine section.
An element whose outer action starts in degree $m>2$ differs, below
that degree, only by $\operatorname{Int}(c)^a$: the only inner special
graded derivation is the degree-two line, as follows from the
centralizer calculation in the proof of G.10 below. Conjugate this
element by a Galois element $t_0$, and multiply by its
$-\chi_\ell(t_0)$ power in the compact pro-$\ell$ image. The
degree-two terms cancel exactly; the degree-$m$ extremal term is
multiplied by
$\psi(t_0)^{m-1}\overline\psi(t_0)-\chi_\ell(t_0)$.
Choose $t_0$ with this nonzero. The operation stays in the image of
$\Delta$, since $\Delta$ is closed and normal, and realizes the
desired line in the genuine graded filtration after rational span.
This argument applies to every packet twist and does not assume
that its degree-two Kummer coefficient is nonzero.
\end{proof}

\begin{statement}{G.10}{Proved}{The full pro-$\ell$ fixed subgroup at every prime}
For every prime $\ell$, including inert primes and $\ell=2$, and
for every pro-$\ell$ cusp-packet section $s_b$ at $O$,
\[
 P_\ell^{s_b(\Delta)}=\overline{\langle c\rangle},
 \qquad P_\ell=\pi_1(U_{\overline K})^{(\ell)}.            \tag{G.12}
\]
In particular this proves the entire pro-$\ell$ assertion in G(b),
including the packet whose degree-two character vanishes.
\end{statement}
\begin{proof}
CM over $K$ makes the homological Galois action abelian at every
prime, so $\Delta$ acts trivially on homology. It fixes inertia
pointwise because $\chi_\ell|_\Delta=1$. All arguments in this
proof that use a Lie expansion are over $\mathbb Q_\ell$, or its
quadratic splitting field; they require no division by $2$ in
$\mathbb Z_2$ and no integral CM eigenspace decomposition.

The degree-one special module is zero. If the genuine degree-two
image of $s_b(\Delta)$ is nonzero, G.3 makes it a nonzero multiple
of $\operatorname{ad}\omega$, whose kernel is $\mathbb Q_\ell\omega$.
Here is a direct proof of that centralizer assertion. The least
associative word of a nonzero homogeneous Lie polynomial is Lyndon.
If it commutes with $\omega$, its least word commutes with $xy$,
and is therefore a power of $xy$. A proper power is not Lyndon.
The remaining degree-two possibility is removed by subtracting a
scalar multiple of $\omega$. Thus the only commuting Lie
polynomials are those multiples.

Suppose instead that the degree-two image vanishes. We obtain a
nonzero metabelian-detected graded derivation using
\citet[Lemma 3.12, Claim 4.14.1, equation 4.14.2, and the proof of
Corollary 4.15]{Nakamura1995}. The metabelian nonvanishing is also stated explicitly in
\citet[Theorem 5.9 and the following paragraph]{Tsunogai1995}.
These statements concern every
rational prime $\ell$, as specified at the start of that paper's
\S1, not only split primes. They provide a degree $m>2$ for which
the actual outer graded image of $G_1$ has nonzero metabelian
projection. The proof of G.2 applies to this outer image at any
prime and transfers its rational span to $\Delta$. There is no
remaining lower genuine inner term: in a symplectic expansion
an inner special derivation has the form $\operatorname{ad}v$
with $[v,\omega]=0$, hence is the central degree-two line just
removed. Induction through the lower degrees therefore identifies
this outer image with the genuine graded image in degree $m$.

Extend scalars to a field splitting the CM action and project
onto one CM character with nonzero metabelian projection.
This projection stays in the arithmetic rational span. One can
verify the required character separation without a density
assumption: use a Frobenius at $5$ when $\ell\ne5$, and one at
$13$ when $\ell=5$. Direct point counts give respectively
$\#E(\mathbb F_5)=8$ and $\#E(\mathbb F_{13})=8$, so the
homological eigenvalues are $-1\pm2i$ and $3\pm2i$.
Their quotients are not roots of unity in $\mathbb Q(i)$.
Thus the eigenvalues on distinct characters $(a,b)$ with
$a+b=m$ are distinct; polynomial spectral projectors in this
Frobenius give exactly the asserted projections. The point
counts are independently recorded in \texttt{g\_arithmetic.py}.
We have obtained a multihomogeneous $D$ with nonzero metabelian
image in the genuine rational graded span of $s_b(\Delta)$.
By G.8 its kernel is the boundary line.

Now let $u\ne1$ be fixed and let $u_d$ be its first nonzero
lower-central term. The leading action error of every actual
fixing automorphism annihilates $u_d$, and so does their rational
span. Either of the two preceding cases therefore forces
$d=2$ and $u_2=a\omega$. The coefficient $a$ is in
$\mathbb Z_\ell$ because the actual boundary has primitive
class in $L_2$. The element $c^{-a}u$ remains fixed and has
zero degree-two term. If it were nonidentity, the same argument
on its first nonzero term would be a contradiction. The lower
central series separates a free pro-$\ell$ group, since every
finite $\ell$-group quotient is nilpotent. This proves (G.12)
in the original group. In particular the extra fixed elements
in the terminal layers of finite $N_r$, exhibited in CK E3,
do not give a compatible non-inertia element in the inverse limit.
\end{proof}

\subsection{Procyclic normalizers: the exact finite stopping level}

\begin{statement}{G.11}{Proved}{A primitive CM eigenline stops at class five}
In the arithmetic pro-$5$ extension, a primitive CM eigenline of
$V$ lifts to a procyclic subgroup of $N_4$ whose arithmetic normalizer
surjects onto $G_K$. No procyclic subgroup of $N_5$ with that primitive
abelian image has arithmetic normalizer with open image in $G_K$.
Thus the largest requested nilpotence class is $r=4$.
\end{statement}
\begin{proof}
Since $4<5$, logarithmic coordinates on the free class-four pro-$5$
group use only $5$-adic units as denominators. Choose CM-equivariant
symplectic coordinates as in G.2, through this class. Compare a Galois
automorphism with its diagonal homological action. Their quotient
fixes the boundary and is IA. The special graded modules in degrees
one and three are zero; degree two is $\mathbb Z_5\operatorname{ad}
\omega$. Consequently this quotient is conjugation by a power of
$c$. Correcting the arithmetic lift by that power makes it preserve
the subgroup $\exp(\mathbb Z_5x)$, proving the normalizer assertion.
The primitive abelian image proves that this subgroup is maximal
procyclic in $N_4$.

For class five, remove inner corrections and work, if necessary,
in the rational nilpotent Lie algebra; this only weakens a necessary
conjugacy condition. By G.7 and the character-separation argument in
G.9, the degree-four outer image of any open subgroup of
$G_K$ contains a nonzero multiple of $T$. If a primitive lift
$h=x+\text{higher terms}$ were carried into its procyclic subgroup
up to conjugacy, an IA Galois element would have exponent $1$ on
that subgroup, by its primitive homological image. At the first
non-inner degree the necessary equation is
\[
                  T(x)\in[x,L_4].                       \tag{G.13}
\]
In bidegree $(3,2)$ the latter space is spanned by $b_{xxxyy}$;
formula (G.7) has the additional coefficient $-1$ at $b_{xxyxy}$.
Thus (G.13) is false. Higher terms of $h$ contribute only in degree
at least six at this comparison, so changing the lift does not
remove the obstruction. Exchanging $x,y$ gives the same conclusion
for the other CM eigenline. The rational character remains nonzero
after passage to an open subgroup, since its image has finite index.
The exact rank check in \texttt{g\_special.py} increases from $1$ to
$2$ when $T(x)$ is adjoined to $[x,L_4]$.
\end{proof}

\begin{statement}{G.12}{Proved}{A normalizer has a boundary leading term at every prime}
Let $C=\overline{\langle u\rangle}\ne1$ be a closed procyclic
subgroup of $P_\ell$. If its arithmetic normalizer has open image
in $G_K$, then the first nonzero lower-central term of $u$ is a
nonzero multiple of $\omega$ in degree two. This holds at every prime.
\end{statement}
\begin{proof}
Let $d$ be the first degree and $\eta$ the conjugation character
on $C\simeq\mathbb Z_\ell$. Its action on the nonzero vector
$u_d\in L_d$ shows that $\eta$ factors through the arithmetic
image, since geometric inner automorphisms are trivial on $L_d$.
At a good Frobenius of norm $q$ its value has complex absolute
value $q^{d/2}$: the graded piece is a quotient of the $d$th
homological tensor power of the elliptic curve.

We compare this with a finite-cover detector. Choose an open normal
$M\subset P_\ell$ with $M\cap C\subset C^\ell$, and put $H=MC$.
The quotient $H/M$ is a nontrivial finite cyclic $\ell$-group
and detects $u$. Hence the image of $u$ in the free
$\mathbb Z_\ell$-module $H^{\rm ab}$ is nonzero. Descend the
geometric cover corresponding to $H$ to a finite extension of $K$
and intersect its arithmetic group with the normalizer of $C$.
The resulting normalizer still has open arithmetic image; after
one further finite extension it surjects onto the new absolute
Galois group. Its action on the detected line in $H^{\rm ab}$
has character $\eta$. The boundary submodule of the homology
of this affine cover has Frobenius weight two and its proper
Jacobian quotient has weight one. At a prime good for both
curves and away from $\ell$, the same value of $\eta$ therefore
has absolute value either $q$ or $q^{1/2}$. It follows that
$d=2$ or $d=1$. This is the finite-cover weight argument underlying
\citet[Lemma 2.1.4 and proof of Theorem 2.1.1]{Nakamura1994};
here it is applied to a detected nonzero homology line before
its weight has been determined.

If $\ell$ is inert in $\mathbb Q(i)$, or $\ell=2$, the representation
$V_\ell$ has no line over $\mathbb Q_\ell$ stable under an open
arithmetic subgroup. Indeed a positive power of a Frobenius at $5$
belongs to that subgroup. Its eigenvalues are powers of
$-1+2i$ and $-1-2i$; their quotient has infinite order and
$i\notin\mathbb Q_\ell$. Thus this power still has irreducible
quadratic characteristic polynomial over $\mathbb Q_\ell$.
So degree one is impossible. At an odd split prime a stable
line is one of the two CM eigenlines. G.11 excludes it at $5$.
The same obstruction works at any other split prime: the
middle degree-four character is nonzero by
\citet[Corollary 4.10]{Ishii2025}, and the two outer degree-four
characters vanish under $[i]$. Its multiple of $T(x)$ fails
membership in $[x,L_4]$ just as in (G.13). This rational
obstruction applies even when the integral homological image
of $u$ is not primitive. Therefore $d=2$ at every prime,
and $L_2=\mathbb Z_\ell\omega$ proves the assertion.
\end{proof}

\begin{statement}{G.13}{Proved}{The pro-$\ell$ arithmetic normalizer criterion}
For every prime $\ell$, every nontrivial closed procyclic subgroup
of $P_\ell$ whose arithmetic
normalizer has open image in $G_K$ is conjugate into
$\overline{\langle c\rangle}$. In particular every maximal such
procyclic subgroup is a cuspidal inertia subgroup.
\end{statement}
\begin{proof}
Write $C=\overline{\langle u\rangle}$ and $N=N_{\Pi^{(\ell)}}(C)$,
where $\Pi^{(\ell)}$ is the geometrically pro-$\ell$ arithmetic group.
The group $P_\ell$ is torsion free, so $C\simeq\mathbb Z_\ell$.
Let $\eta:N\to\mathbb Z_\ell^\times$ be its conjugation character.
By G.12 the first nonzero term of $u$ is $a\omega$ in $L_2$,
with $a\ne0$. For every $n\in N$, projecting its conjugation
action onto this free rank-one module gives
\[
             a\eta(n)\omega=a\chi_\ell(p(n))\omega.
\]
Cancellation in $\mathbb Z_\ell$ proves $\eta=\chi_\ell\circ p$
on all of $N$, including its geometric kernel. Thus $N$ is a
cyclotomic normalizer in the exact sense of
\citet[\S2.1, Definition preceding Theorem 2.1.1]{Nakamura1994}.
His Theorem 2.1.1 applies to every full class of finite groups,
including the finite $\ell$-groups, and to every affine hyperbolic
type, including $(1,1)$. It puts $u$ in an actual conjugate of
cuspidal inertia in $P_\ell$. This proves the assertion without an
assumption about integral conjugacy in the Malcev completion.
For clarity, the theorem's decisive finite-cover statement is its
Lemma 2.1.4: a noncuspidal element survives modulo the boundary
part of the abelianization of a suitable finite cover. A
cyclotomic normalizer would give a weight-two line there, while
the remaining Jacobian quotient has weight one. The resulting
contradiction concerns the original pro-$\ell$ group. G.12 has
verified its cyclotomic hypothesis at every prime.
\end{proof}

\subsection{Computation through degree fourteen and its arithmetic bounds}

\begin{statement}{G.14}{Proved}{The ambient modules and bracket images}
The following entries are exact. The second row is
$\dim\mathfrak s_m$; the third is the dimension of its $[i]$-fixed
part. Brackets are brackets of all strictly lower positive special
degrees; the even-only row is also given because arithmetic odd
degrees vanish.
\[
\begin{array}{c|rrrrrrr}
m&2&4&6&8&10&12&14\\\hline
\mathfrak s_m&1&3&6&13&37&99&284\\
\mathfrak s_m^{[i]}&1&1&4&5&21&45&150\\
{}[\mathfrak s_{>0},\mathfrak s_{>0}]_m&0&0&0&3&15&47&153\\
\text{even lower brackets}&0&0&0&3&15&47&147\\
{}[\mathfrak s_{\rm even}^{[i]},\mathfrak s_{\rm even}^{[i]}]_m
 &0&0&0&0&4&11&41
\end{array}                                             \tag{G.15}
\]
All CM character multiplicities are supplied in
\texttt{code/g\_special\_results.json}; every weight $(a,b)$ there
means $\psi^a\overline\psi^b$. The free Lie ranks through degree
fourteen are
\[
2,1,2,3,6,9,18,30,56,99,186,335,630,1161.
\tag{G.16}
\]
\end{statement}
\begin{proof}
For $L_{r;(a,b)}$, with $a+b=r$, the integral Witt rank is
\[
 \ell(a,b)=\frac1{a+b}\sum_{d\mid\gcd(a,b)}
      \mu(d)\binom{(a+b)/d}{a/d}.
\]
Equation (G.2) gives the special character multiplicity
$\ell(a+1,b)+\ell(a,b+1)-\ell(a+1,b+1)$ in degree $a+b$.
Keeping $a-b\equiv0\pmod4$ gives the fixed part. The program
enumerates all binary Lyndon words through degree sixteen,
checks the triangular leading coefficient $1$ of every standard
bracketing through degree fourteen, and checks the Witt counts.
It constructs primitive integral representatives of rational
nullspace bases of (G.2), by character, through degree ten, and
computes derivation brackets with exact rational elimination.
Only degree-seven odd derivations can contribute an odd--odd
bracket through degree fourteen. Adding their six independent
degree-fourteen brackets changes $147$ to $153$; this check is in
\texttt{g\_kernel.py}. The basis vectors used for rational rank
calculations are not claimed to be saturated integral bases.
The integral module ranks instead follow from the integral exact
sequence (G.2). This distinction avoids drawing an integral
conclusion from a rational nullspace computation.
\end{proof}

\begin{statement}{G.15}{Known}{Exact rational arithmetic dimensions through fourteen}
For the outer graded Galois image at $(K,\ell)=(\mathbb Q(i),5)$,
the rational dimensions in even degrees are
\[
\begin{array}{c|rrrrrrr}
m&2&4&6&8&10&12&14\\\hline
\text{new metabelian-detected generators}&0&1&3&3&5&5&7\\
\text{rational outer Galois image}&0&1&3&3&8&11&24
\end{array}                                             \tag{G.17}
\]
This computation uses the opened 2026 preprint
\citet[Theorem 1.2 and Lemma 1.5]{Ishii2026Graded}; it is not
a claim of freeness in unbounded degree. All CM character
multiplicities are in \texttt{g\_arithmetic\_results.json}.
The genuine standard-tangent image additionally has
the central degree-two line of G.4.
\end{statement}
\begin{proof}
The cited theorem gives one generator for each pair
$(a,b)\in\mathbb Z_{>0}^2$, $a\equiv b\pmod4$,
$(a,b)\ne(1,1)$, assuming the indicated $H^2$-vanishing.
We verify the arithmetic criterion in its Lemma 1.5 explicitly.
The ray class field $K(5)$ is $\mathbb Q(\zeta_{20})$:
both extensions of $K$ have degree $4$ and conductor dividing
$(5)$. Its discriminant is $4{,}000{,}000$, and the degree-eight
Minkowski bound is strictly between $12$ and $13$. The residue
degrees at $2,3,7,11$ are $4,4,4,2$, so the only prime ideals
of norm at most that bound are the two primes over $5$.
The element $\zeta_{20}^3-\zeta_{20}+1$ has norm $5$;
its conjugate generates the other prime. Hence the class number
is $1$. For $K(25)$, the ray class group is
$(\mathbb Z/25)^\times{}^2/\mu_4$, of order $100$.
At $\mathfrak p=(2+i)$ the inertia group has order $20$;
in the quotient of order $5$, Frobenius is represented by
$2+i\equiv9\pmod{\overline{\mathfrak p}^{\,2}}$, using
$i\equiv7\pmod{25}$. Its coset generates
$(\mathbb Z/25)^\times/\mu_4$. Thus the decomposition
group is the whole ray group. The conjugate prime behaves
identically, and there are exactly two primes over $5$ in
$K(25)$. All numerical identities, norms, and cosets are
checked in \texttt{g\_arithmetic.py}. These establish the
precise hypotheses of the cited lemma.

Each arithmetic generator is detected by a distinct nonzero
metabelian character. Brackets have zero metabelian character,
so the new generators are independent modulo all lower
brackets. Write $S_m=\mathfrak s_m^{[i]}$. The exact ranks
in G.14 prove injectivity of the entire ambient maps
\[
\begin{aligned}
 S_4\otimes S_6&\longrightarrow S_{10},& 1\cdot4&=4,\\
 (S_4\otimes S_8)\oplus\bigwedge^2S_6
   &\longrightarrow S_{12},& 1\cdot5+6&=11,\\
 (S_4\otimes S_{10})\oplus(S_6\otimes S_8)
   &\longrightarrow S_{14},& 1\cdot21+4\cdot5&=41.
\end{aligned}
\]
Restrict these injective maps to the actual arithmetic
subspaces; no choices of their unknown multiplicity-space
coordinates affect injectivity. Through degree eight the
image dimensions are the generator counts $1,3,3$. In
degree ten the brackets add $1\cdot3$ to $5$, giving $8$.
In degree twelve they add $1\cdot3+\binom32$ to $5$, giving
$11$. In degree fourteen they add $1\cdot8+3\cdot3$ to
$7$, giving $24$. Every bracket is included: the positive
arithmetic degrees below four vanish, and the degree-two
genuine line is central. Character addition gives the
multigraded table in the output. This proves equality
through fourteen without assuming the preprint's
unbounded freeness conjecture.
\end{proof}


For comparison with B.18 and the supplied LM3b A.7--A.8,
A.12--A.16, subtracting the central degree-two line from (G.15)
gives the outer ambient dimensions
$0,3,6,13,37,99,284$. The universal geometric elliptic image in
those degrees is $0,3,5,10,24,50,109$, as computed in B.18 and
its degree-fourteen check. It is a different image from the
specialized CM arithmetic row (G.17). In particular the universal
$\mathrm{SL}_2$-orbit in degree four has dimension three, whereas
the CM arithmetic image has dimension one. The relations of
\citet[Theorem 1]{Pollack2009} concern the epsilon-generated
geometric algebra; they do not identify its specialization with
this arithmetic image.

\begin{statement}{G.21}{Proved}{The actual degree-six middle line}
Use the symplectic normalization of the projector eigenbasis of G.7:
if $e_+,e_-$ have columns $(1,-i_5)/2,(1,i_5)/2$ in the original
symplectic basis, put $d=\det(e_+,e_-)=i_5/2$ and
$x=e_+,\ y=d^{-1}e_-$. Thus $(x,y)$ is symplectic;
the symbol $d$, with $d^2=-1/4$, denotes the determinant of the
\emph{unnormalized} projector basis.
Let $F(x)=y,F(y)=0$, and let $Z_3$ be the special derivation
\[
\begin{aligned}
 Z_3(x)&=\tfrac14 b_{xxxyxyy}+\tfrac14 b_{xxxyyxy}
                +\tfrac1{12}b_{xxyxyxy},\\
 Z_3(y)&=-\tfrac14 b_{xxyxyyy}-\tfrac14 b_{xxyyxyy}
          -\tfrac14 b_{xxyyyxy}-\tfrac1{12}b_{xyxyxyy}.
\end{aligned}                                                    \tag{G.21}
\]
Then $Z_3$ is $\mathrm{SL}_2$-invariant and has zero metabelian image.
The degree-six arithmetic subspace of character $(3,3)$ is the line
\[
 \mathbb Q_5 V_6,\qquad
 V_6=\frac{[F,[F,\epsilon_6]]-192Z_3}{6}.                \tag{G.22}
\]
More precisely, on the sixth Johnson subgroup the $(3,3)$ projection is
$\chi_3 V_6/1536$, where $\chi_3$ is the raw cyclotomic Soul\'e
character of \citet[Definition 1.5]{IshiiGoto2026}. Thus the actual
rational degree-six space is spanned by $\epsilon_6$, $V_6$, and
$[F,[F,[F,[F,\epsilon_6]]]]$. These are the characters $(5,1)$,
$(3,3)$ and $(1,5)$, respectively. This assertion holds for every
cusp packet; no equality of integral arithmetic lattices is asserted.
\end{statement}
\begin{proof}
The exact special-kernel calculation gives the canonical
$\mathrm{GL}_2$-decomposition
\[
 \mathfrak s_6=(\operatorname{Sym}^4 V\otimes\det V)
                 \oplus(\det V)^3.
\]
The metabelian projection is an isomorphism on the first summand and
zero on the second. The latter is spanned by $Z_3$. Direct substitution
in (G.21) proves $Z_3(\omega)=0$, $[F,Z_3]=0$, and
$[x\partial_y,Z_3]=0$. Its metabelian image is zero. Also
\[
 h([F,[F,\epsilon_6]])=12X_0^2Y_0^2.                   \tag{G.23}
\]
These are exact free-Lie calculations, reproduced by
\texttt{g\_embedding.py}.

We justify the arithmetic coefficient of the invariant summand,
rather than treating a Tate-fibre formula as a CM-fibre formula.
The universal mixed elliptic motive and its lisse \(5\)-adic realization
have a natural quotient given by constant mixed Tate objects.
See \citet[Definition 6.1, the exact sequence in \S7.1,
Proposition 20.1, Corollary 20.2, and \S22]{HainMatsumoto2020}.
In weight six, the universal unipotent Lie algebra has exactly the
geometric $\operatorname{Sym}^4V\otimes\det V$ generator and the
constant $\mathbb Q_5(3)$ generator. There are no brackets of lower
noncentral generators in this weight, and the degree-two generator
is central. Its graded monodromy map is injective in weight six:
the geometric generator is detected metabelianly, and the explicit
constant generator (G.21) is nonzero in the complementary invariant
line. Consequently the coefficient on the quotient
$\mathfrak s_6/(\operatorname{Sym}^4V\otimes\det V)$ at any smooth
elliptic fibre is the coefficient of this same constant mixed Tate
representation. Pullback along the point $E/K$ restricts that
representation from $G_{\mathbb Q}$ to $G_K$. Transport of a
symplectic homology basis changes neither the invariant vector nor
the quotient coefficient. In degrees at most six, changes of a
symplectic expansion cause no further ambiguity: degree-one and
degree-three special derivations vanish, and the degree-two one is
central. This is a functorial statement about the \emph{quotient};
the geometric coefficient at the CM fibre must still be computed.

For the exact normalization of the constant quotient, it is essential
to distinguish the two filtrations. Write $W$ for the ordinary
elliptic weight, so $x$ and $y$ have $W$-weight $-1$ and derivation
degree $r$ means $W$-weight $-r$. At the first-order Tate tangent
the relative monodromy filtration $M$ gives $x$ weight $0$ and the
vanishing cycle $y$ weight $-2$. The mixed Tate Lie algebra
$\mathfrak k$ has a one-dimensional $\operatorname{Gr}^M_{-6}$,
and no negative $M$-weights greater than $-6$; see
\citet[\S3 and Proposition 20.1]{HainMatsumoto2020}.
The Tate section $\phi_E:\mathfrak k\longrightarrow
\operatorname{Der}\mathfrak p(E_{\mathrm{Tate}}')$ respects $M$
but does not respect $W$. The precise comparison is
\citet[Proposition 29.1 and Corollary 29.2]{HainMatsumoto2020}:
all $W$-components of an $M$-weight $-6$ generator except the
$W$-weight $-6$ component are geometric, and its $W$-weight $-6$
component represents the same constant quotient generator.

Take its explicitly normalized Tate lift $\Sigma_3$ of
\citet[(5.41), (5.46), and (5.47)]{DorigoniEtAl2026}. Its components
in degrees four and six are
\[
        (\Sigma_3)_4=-\tfrac12[F,[F,\epsilon_4]],
        \qquad (\Sigma_3)_6=Z_3.
\]
Both components have $M$-weight $-6$: their characters in the Tate
basis are $(1,3)$ and $(3,3)$, respectively. These formulas can be
checked without a period comparison.
Put $s_{12}=[y,x]$, $s_{01}=-y-\sum_{n\ge1}B_n
\operatorname{ad}_x^n(y)/n!$, and $g_3(u,v)=[u-v,[u,v]]$.
They satisfy the defining equation
$\Sigma_3(s_{01})=[s_{01},g_3(s_{12},-s_{01})]$
through length seven, as the supplied exact calculation verifies.
If $\beta(\sigma)$ is the coefficient of this normalized generator
in the $M$-weight $-6$ part of the mixed Tate Galois action, then
the entire $M$-weight $-6$ Tate derivation is
$\beta(\sigma)\Sigma_3$. There is no independent geometric
coefficient in this assertion: it concerns the image of the
specified Tate section $\phi_E$, whose source in that weight is
one-dimensional. Nor is there a bracket of lower mixed Tate
generators in this weight. Its ordinary-degree-four projection is
therefore $\beta(\sigma)(\Sigma_3)_4$.
Since $h((\Sigma_3)_4)=-Y_0^2$, the degree-two term of Nakamura's
invariant is $\beta(\sigma)Y_0^2$. The Tate formula
\citet[Theorem 5.2.6]{Nakamura2013Eisenstein} is
\[
       E_{\sigma,2}=\frac{\chi_3(\sigma)}{2(1-5^2)}Y_0^2
                     =-\frac{\chi_3(\sigma)}{48}Y_0^2.
\]
Thus $\beta=-\chi_3/48$. Proposition 29.1 identifies the
ordinary-degree-six projection modulo geometric derivations with
$\beta Z_3$; there are no degree-six geometric $\mathrm{SL}_2$
invariants, so $Z_3$ is the unique invariant representative.
Only this constant quotient coefficient is transported to the CM
fibre. The geometric ordinary-degree-four term at the Tate
tangent is used to measure $\beta$, and is not transported to
the CM fibre.
There is an apparent sign error in the second formula of
\citep[(5.53)]{DorigoniEtAl2026}: the defining equation and (5.46)
give $+(b_{xyyyy})$ as the value on $y$. Using the printed minus
sign would give the nonzero boundary error $-2b_{xxyyyy}$.
Neither this proof nor the code uses that incompatible sign.

It remains to compute the CM geometric coefficient. The raw theta
character satisfies
\[
 \chi^{\rm ell}_{(3,3)}
 =-3\cdot4\frac{B_{3,\chi_{-4}}}{3}
     \frac{1-5^4}{1-5^2}\chi_3=-156\chi_3,
 \qquad B_{3,\chi_{-4}}=\tfrac32,
\]
by \citet[Theorem 1.7]{IshiiGoto2026}. As in G.7, both the
determinant correction and Nakamura's published factor-twelve erratum
are necessary. Thus
\[
       12\kappa_{22}=d^2\chi^{\rm ell}_{(3,3)}=39\chi_3,
       \qquad \kappa_{22}=\tfrac{13}{4}\chi_3.
\]
The coefficient of $X_0^2Y_0^2$ in the invariant $\alpha=-h$ is
\[
 \frac{\kappa_{22}}{(1-5^4)2!2!}=-\frac{\chi_3}{768}.
\]
Using (G.23), its geometric component is
$\chi_3[F,[F,\epsilon_6]]/9216$. Adding the constant quotient
component $-\chi_3Z_3/48$ proves (G.22) and its stated coefficient.
Nonvanishing in the sixth Johnson quotient follows from
\citet[Corollary 4.10]{Ishii2025} and the distinct-weight argument
of \citet[Claim 4.14.1 and (4.14.2)]{Nakamura1995}; the central
degree-two character has a different weight and can be removed by
the same argument. Alternatively the exact arithmetic dimensions
in G.15 give this nonvanishing. The two extreme character spaces
are already one-dimensional in the ambient module, and their
nonzero arithmetic images were established in G.15. Finally
packet changes add only central inner inertia, so they have no
effect on this degree-six outer space.
\end{proof}

\begin{statement}{G.16}{Open}{The remaining embedded arithmetic image}
G.15 determines all rational character dimensions requested
through degree fourteen, and G.21 determines the actual rational
embedding through degree six. These do not determine the actual
$\mathbb Z_5$-lattice and the coordinates of every arithmetic
generator in each repeated-character special-derivation
space. Those integral and coordinate data remain an
undecided portion of G(e).
\end{statement}
\begin{proof}[Attack and remaining step]
The character decomposition in (G.15) cuts the ambient spaces
down substantially, and the ambient bracket injections
give the exact rational dimensions in (G.17). The nonzero
metabelian characters detect the abelianized generators;
they do not determine their components
in the kernel of the metabelian comparison. G.21 resolves the
first ambiguity, in degree six and character $(3,3)$, using the
universal constant mixed Tate quotient in addition to the
metabelian character. The first remaining multiplicity-space
placement in this return is degree eight, character $(4,4)$.
There the ambient special space has dimension three; the
universal geometric subspace still has a two-dimensional middle
character space, whereas the arithmetic image has dimension one.
The Tate-quotient argument of G.21 supplies no even Tate-weight
generator to determine the remaining ratio. We can make that ratio
and the failure of the finite tests entirely explicit. Set
\[
 A_8=\frac{\operatorname{ad}_F^3(\epsilon_8)}{120},
 \qquad
 B_8=\frac{[\epsilon_4,\operatorname{ad}_F^2(\epsilon_4)]}{2}.
 \tag{G.24}
\]
Here $B_8$ is an integral primitive nonzero vector. The exact
calculation in \texttt{g\_degree8.py} gives
$h(A_8)=X_0^3Y_0^3$ and $h(B_8)=0$. They form a basis of the
degree-eight, character-$(4,4)$ universal geometric space: its
two summands are the middle lines of the $\epsilon_8$ module and
of $\Lambda^2\langle\epsilon_4,F\epsilon_4,F^2\epsilon_4\rangle$.
There is no constant mixed Tate generator in this ordinary
weight, and brackets with the degree-two central generator are
zero; this follows from \citet[Proposition 20.1 and
\S22]{HainMatsumoto2020}, or the exact universal geometric table
in Part B. Consequently, for a uniquely determined but presently
unevaluated $\lambda\in\mathbb Q_5$, the actual middle projection
on the eighth Johnson subgroup is
\[
 \frac{d^3\chi^{(4)}_4}{267840}\,(A_8+\lambda B_8).
 \tag{G.25}
\]
To verify the scalar, \citet[Theorem 1.7]{IshiiGoto2026} gives
$\chi^{\rm ell}_{(4,4)}=126\chi^{(4)}_4/5$.
The same published factor-twelve correction as in G.7 gives
$\kappa_{33}=d^3(21/10)\chi^{(4)}_4$, whence
$h=-\kappa_{33}X_0^3Y_0^3/((1-5^6)(3!)^2)$ has the coefficient
in (G.25). The one-dimensional actual image in G.15 and its
nonzero metabelian projection prove existence and uniqueness of
$\lambda$; they do not give its value.

There is a precise algebraic obstruction to determining it from
the tests already used. Both $A_8$ and $B_8$ are special, are
fixed by the order-four CM action, and acquire sign $-1$ under
the interchange $x\leftrightarrow y$. Thus every
$A_8+\lambda B_8$ has the same prescribed metabelian coefficient
and those symmetries. The three ambient bracket injections of
G.14 show that replacing this one arithmetic generator by any
such vector, while keeping the other chosen generators, leaves
all the rational dimension and lower-bracket tests through
degree fourteen unchanged. The dimensions remain
$0,1,3,3,8,11,24$. This is an algebraic family of candidate
truncated subalgebras, not a claim that every parameter is
realized by Galois. Although the actual degree-four arithmetic
space is the single line $\mathbb Q_5T$ and hence has zero
self-bracket, $B_8$ is a bracket of the larger universal
geometric space. Its coefficient in a new degree-eight
arithmetic generator is not forced to vanish by
$[T,T]=0$. The missing arithmetic datum is exactly this first
scalar $\lambda$, followed by the higher specialized generator
vectors and their integral indices.
The rational rank calculations do not supply those indices.

We also tested the exact higher Tate formula, rather than merely
stopping at the absence of an even generator. Equations (5.46)
and (7.22) of \citet{DorigoniEtAl2026} give
\[
       (\Sigma_3)_8=\frac{[\epsilon_4,F\epsilon_4]}{480},
       \qquad [F,(\Sigma_3)_8]=\frac{B_8}{240}.
       \tag{G.26}
\]
The program verifies the defining lift (5.41) through word length
nine, including this term. In the Tate basis its character is
$(5,3)$, so its relative $M$-weight is $-6$, although its
ordinary $W$-weight is $-8$. It is entirely geometric and maps
to zero in the constant quotient; see that paper's
Theorem 5.4.1(v), or
\citet[Proposition 29.1]{HainMatsumoto2020}.
Thus (G.26) determines a geometric component of the Tate lift
of $\chi_3$, not the CM specialization coefficient of
$\chi^{(4)}_4$ in (G.25). The geometric transport missing from
the degree-eight calculation is exactly what the quotient
argument of G.21 was able to avoid in degree six.

Nor can one impose the Tate recursion's extra condition
$[\Sigma_w,N]=0$ from
\citet[Theorem 5.4.1(vii), equation (5.43)]{DorigoniEtAl2026}
on the CM arithmetic generator. The ordinary-degree-zero term
of $N$ is $-F$, so such a condition on a derivation beginning
in degree eight would require
$[F,A_8+\lambda B_8]=0$. The exact calculation gives rank two
for the pair $[F,A_8],[F,B_8]$, ruling this out for every
$\lambda$. This is a failed transfer of a condition specific
to the Tate tangent, not an arithmetic equation selecting
the parameter. The opened all-order formula therefore does
not close G(e); a CM geometric specialization calculation is
still required.
We also applied the newer local reciprocity formula of
\citet[Theorem 1.1, Corollary 5.5 and Remark 5.6]{GotoIshii2026Coleman}.
In that paper's normalization, specializing to $K=\Q(i)$, $p=5$
and $(i,j)=(4,4)$ gives
\[
 \chi^{\rm ell,new}_{(4,4)}(\operatorname{rec}_\infty\beta)
 =-6048\,L_p(-3,-3,\omega_{(4,4)}^{-1})
          \phi^{\rm CW}_{(4,4)}(\beta).
\]
Indeed the scalar is
$48(1-5^3)(1-2/(1-5^3))=-6048$.
The reciprocity map is from the specified norm-compatible semilocal
principal units to the product of local abelian Galois groups; it
does not parametrize arbitrary elements of the eighth Johnson subgroup.
This formula determines a restriction of the metabelian character,
and hence annihilates $B_8$. It therefore supplies no equation for
$\lambda$. Integral image indices would additionally require the
image lattice of the Coates--Wiles map and the comparison with that
Johnson subgroup. The normalization correction in Remark 5.6
concerns the older Coleman--Ihara corollary, not Theorem 1.7 used
in (G.25).
The genuinely nonmetabelian follow-up is G.23--G.24. Enriquez's
universal elliptic kernel contains the entire one-parameter
family after exponentiation, and finite-group averaging transports
it to any prescribed exact finite CM lift without changing its
leading vector. Thus imposing all universal algebraic equations,
rather than only their finite truncations, does not select
$\lambda$. G.24 instead supplies an integral second coordinate
that does detect $B_8$, with an exact finite Magnus quotient
on which to measure it. The missing datum is the value of that
coordinate on an actual eighth-Johnson Galois element with
nonzero first coordinate; no such measurement is supplied here.

We also opened a source that does address depth two rather than
the metabelian quotient. \citet[Theorems 11 and 16]{Malkin2020CM}
proves second-shuffle and Bianchi-complex results for symmetric
Hodge correlators on CM elliptic curves punctured at torsion
points. Its \S2.2.3, Lemma 3 and the following paragraph
(printed pp.~12--13) explicitly make the passage to the
$\ell$-adic realization conditional on the motivic formalism
and the stated regulator-injectivity hypothesis. The averaged
basepoint construction is also different from the specified
once-punctured tangent; see its \S4.3.2 remark after Theorem 15.
Consequently that theorem is not an evaluated arithmetic value
of the coordinate in G.24. Neither an unproved realization
comparison nor a Hodge-period guess is used to assign $\lambda$.

As a separate attack on the exceptional fixed
packet, \texttt{g\_kernel.py} computes the kernel of $T$ on
$L_d$ for $1\le d\le10$: it is one-dimensional only at $d=2$.
Those finite tests were superseded for fixed subgroups by the
unbounded argument G.8--G.10, but are retained as reproducible
evidence.
\end{proof}

\begin{statement}{G.23}{Proved}{Universal elliptic equations and exact finite CM transport leave the parameter free}
Let $k=\mathbb Q_5$, with the chosen embedding of $\mathbb Q(i)$,
and use $A_8,B_8$ from (G.24). For every $\lambda\in k$,
\[
             t\longmapsto\exp\bigl(t(A_8+\lambda B_8)\bigr)
\]
is a one-parameter subgroup of Enriquez's graded universal elliptic
Grothendieck--Teichm\"uller kernel $R_{\rm ell}^{\rm gr}$.
It commutes with $J(x)=ix$, $J(y)=-iy$; the interchange
$S(x)=y$, $S(y)=x$ acts on it by $t\mapsto-t$.
The full diagonal CM torus normalizes it through
$(a,b)\mapsto a^4b^4$.

Thus all defining universal elliptic group equations, together
with these character and finite-symmetry conditions, admit every
$\lambda$ to all orders. The same assertion about the leading
degree-eight vector holds for any prescribed lift of the finite
group $\langle J\rangle$, or $\langle J,S\rangle$, in the
corresponding $\mathrm{GL}_2$-extension of the positive
prounipotent kernel. The determinant-$-1$ element $S$ is taken
in this extension, not in the $\mathrm{SL}_2$ quotient of
$R_{\rm ell}^{\rm gr}$. These are universal group
elements, not asserted elements of the actual arithmetic subgroup.
\end{statement}
\begin{proof}
In \citet[Proposition 4.19, equation (48)]{Enriquez2014},
$\delta_{2n}$ is our $\epsilon_{2n+2}$: pairing the two halves
of his sum for its value on $y$ gives G.8. His bracket is
written in the opposite derivation algebra; changing its sign
does not change the generated Lie subalgebra or the membership
argument. The same section defines
$\mathfrak b_3=\langle\mathfrak{sl}_2,\delta_{2n}\rangle$ and
gives the subgroup inclusion
\[
 \exp(\widehat{\mathfrak b}_3^+)\rtimes\mathrm{SL}_2
                      \subset R_{\rm ell}^{\rm gr}.
\]
Both $A_8$, an iterated $\mathfrak{sl}_2$ bracket with
$\epsilon_8$, and $B_8$, a bracket of two members of the
$\epsilon_4$ module, belong to $\mathfrak b_3$.
Exponentiation converges in the positive-degree filtration.
This is actual all-order group membership, not only a tangent
vector solving finitely many equations. The elliptic-associator
isomorphisms of \citet[\S4.8 and Proposition 4.20]{Enriquez2014}
also transport these elements to the ungraded universal group,
with the canonical associated graded maps. No generation
conjecture for $\mathfrak b_3$ is used.

Both vectors have character $(4,4)$ and change sign under $S$.
Their exact associative expansions are checked in
\texttt{g\_degree8\_nonmet.py}. Hence $J$ fixes them and the
diagonal torus scales them by $a^4b^4$. The homogeneous
$\mathrm{SL}_2$-stable algebra $\widehat{\mathfrak b}_3^+$
is also $\mathrm{GL}_2$-stable, so these determinant-$-1$
and diagonal actions are defined in
$\exp(\widehat{\mathfrak b}_3^+)\rtimes\mathrm{GL}_2$.
Exponentiation proves the asserted identities.

Here is the finite-lift assertion with its convergence argument.
Let $H$ be either finite group just specified and $U$ the
positive-degree prounipotent kernel. Every section
$H\to U\rtimes H$ is conjugate by $U(k)$ to the linear section.
On a successive central vector layer, a cocycle satisfies
$a_{gh}=a_g+g a_h$. For $b=|H|^{-1}\sum_h a_h$ one has
$a_g=b-gb$, so a conjugation removes that layer.
Repeating in increasing degrees gives corrections tending to
the identity, whose product converges in $U(k)$ and conjugates
the entire section. Conjugating the one-parameter subgroups by
the same element supplies the prescribed finite action.
Their leading degree-eight terms remain unchanged, because
the conjugator has positive degree.

The analytic elliptic associator at the CM point, its comparison
path with the Tate tangent, and its arithmetic realization carry
additional information. This proof does not specify those data.
It rules out selecting $\lambda$ from the stated universal
algebraic conditions, even with exact finite CM lifts, and
does not assert arithmetic realization of the candidate family.
\end{proof}

\begin{statement}{G.24}{Proved}{A saturated ambient plane and two finite nonmetabelian coordinates}
Let $[w]P$ denote an associative-word coefficient, and put
\[
 \begin{aligned}
 w_0&=x^5y^4,& w_1&=x^4yxy^3,\\
 \alpha(D)&=-[w_0]D(x),&
 \nu(D)&=[w_1]D(x)+4[w_0]D(x),\\
 C_8&=A_8+\frac{31}{20}B_8.
 \end{aligned}
\]
For $H=\mathbb Z_5x\oplus\mathbb Z_5y$, the plane
$W_8=\mathbb Q_5A_8\oplus\mathbb Q_5B_8$ has saturated ambient lattice
\[
 W_8\cap\operatorname{Hom}_{\mathbb Z_5}(H,L_9H)
               =\mathbb Z_5C_8\oplus\mathbb Z_5B_8.
\]
The coordinate matrix is
\[
 \begin{array}{c|rrr}
       &A_8&B_8&C_8\\\hline
 \alpha&1&0&1\\
 \nu&-31/20&1&0
 \end{array}.
\]
For an actual Galois element $\sigma$ in the eighth Johnson
subgroup these functionals isolate its character-$(4,4)$ part,
and
\[
 \begin{aligned}
 \alpha(D_8(\sigma))&=\frac{d^3\chi^{(4)}_4(\sigma)}{267840},\\
 \nu(D_8(\sigma))&=\alpha(D_8(\sigma))
                         \left(\lambda-\frac{31}{20}\right).
 \end{aligned}
\]
Thus any such element with nonzero $\alpha$ would determine
$\lambda=31/20+\nu/\alpha$.

For each $r\ge1$, both coordinates modulo $5^r$ are read from
the action on a characteristic finite quotient of $P_5$, namely
its image in the units of
\[
 (\mathbb Z/5^r)\langle X,Y\rangle/(X,Y)^{10},
 \qquad a\longmapsto1+X,\quad b\longmapsto1+Y,
\]
where $a,b$ lift the chosen integral homology basis.
This measurement uses no logarithm or division by $5$.
The lattice above is ambient; neither the arithmetic image
lattice nor the arithmetic value of $\nu$ is determined here.
\end{statement}
\begin{proof}
The exact tensor calculation gives
\[
 \begin{array}{c|rrr}
       &A_8(x)&B_8(x)&C_8(x)\\\hline
 [w_0]&-1&0&-1\\
 [w_1]&49/20&1&4
 \end{array}.
\]
Both $C_8$ and $B_8$ annihilate $[x,y]$. All their standard
Lyndon coefficients have denominators in $\{1,2,4\}$; the full
Lyndon vectors are in \texttt{code/g\_degree8\_nonmet.json}
and \texttt{code/g\_degree8\_results.json}. Hence both are integral
Lie derivations at $5$. The integral functionals $\alpha,\nu$
have the identity matrix on the pair $(C_8,B_8)$. Every integral
$D\in W_8$ therefore has
\[
                  D=\alpha(D)C_8+\nu(D)B_8
\]
with coefficients in $\mathbb Z_5$. This proves saturation;
it makes no assertion of arithmetic surjectivity.

Each selected word in $D(x)$ has five $x$'s and four $y$'s,
so both functionals kill all other homogeneous CM characters.
Formula (G.25) now gives the two arithmetic identities, with
its corrected theta-character normalization unchanged.
The second coordinate is nonmetabelian since $\nu(B_8)=1$
and the metabelian character kills $B_8$. It does not vanish
on our chosen lift $A_8$; the explicit $31/20$ corrects that
fact.

Choose lifts $a,b$ of $x,y$ in $P_5$. They form a basis modulo
the Frattini subgroup and are free pro-$5$ generators.
The completed group algebra identifies with the completed
noncommutative power-series algebra by $a\mapsto1+X$,
$b\mapsto1+Y$. The inverse maps follow directly from the free
pro-$5$ universal property and the substitutions
$X\mapsto a-1$, $Y\mapsto b-1$.
The ideal generated by $5^r$ and the tenth power of the
augmentation ideal is characteristic. Its quotient algebra
is finite, and the group $1+(X,Y)$ in it is a finite $5$-group:
its successive augmentation quotients are additive finite
$5$-groups. This proves the claimed characteristic finite
group quotient and the descent of the Galois action.

Since $\sigma$ is in the eighth Johnson subgroup,
$\sigma(a)a^{-1}\in\gamma_9P_5$. The length-nine part of its
Magnus series is $D_8(\sigma)(X)$; all positive lower parts
vanish. Indeed the integral associated graded of a free
pro-$5$ group is its free Lie algebra, and leading group
commutators expand as the corresponding tensor Lie brackets.
Applying $-[X^5Y^4]$ and $[X^4YXY^3]+4[X^5Y^4]$ to that
series gives the two coordinates modulo $5^r$.
Changing a logarithmic expansion cannot change this leading
Johnson term.

The code also verifies detection using actual group words.
At modulus $125$ and length nine, products of the standard
Lyndon group commutators with the coefficients of $C_8$ and
$C_8+B_8$ give central corrections to the two generators.
Their denominators are powers of $2$, so these powers of
commutators are well-defined. The program recursively expands
each group commutator, checks its leading Lie polynomial, and
checks that the resulting substitutions have inverses and
preserve the peripheral commutator. Their measured pairs
$(\alpha,\nu)$ are $(1,0)$ and $(1,1)$ modulo $125$.
These are ambient finite-quotient automorphisms, not claimed
Galois elements. G.15 ensures a nonzero arithmetic first
coordinate in the rational graded image, but its integral
image ideal and the corresponding second coordinate remain
unevaluated. The value of $\lambda$ therefore remains Open.
\end{proof}



\subsection{The first failure of abelian class-field factorization}

\begin{statement}{G.17}{Proved}{All nilpotence levels for the specified tangent}
For the standard tangent $t=X/Y$, the genuine $G_K$-action on
$N_r$ factors through $G_K^{\rm ab}$ exactly for $r=1,2$.
For $r=3$ the first obstruction on $\Delta$ is
\[
 \delta\longmapsto\kappa_{1+i}(\delta)\operatorname{ad}\omega,
 \tag{G.18}
\]
and its fixed field over $K^{\rm ab}$ is
\[
       K^{\rm ab}\bigl((1+i)^{1/5^\infty}\bigr).
 \tag{G.19}
\]
For the exceptional packet class of G.5 the action factors
through $G_K^{\rm ab}$ through $N_4$ and first fails at $N_5$;
the first field is the field cut out by
$\chi^{(4)}_2|_\Delta$, equivalently by
$\chi^{\rm ell}_{(2,2)}|_\Delta$.
\end{statement}
\begin{proof}
CM makes the action on $N_1$ abelian. The genuine order-two
lift $J^2$ and the invertibility of $2$ eliminate the central
shear in $N_2$, as in CK E1; equivalently special degree one
is zero. G.4 gives the first nontrivial $\Delta$-action on
generators modulo $\gamma_4P$, proving (G.18). Kummer theory
gives precisely (G.19), since $K^{\rm ab}$ contains all
$5$-power roots of unity. Every higher $N_r$ maps
equivariantly onto $N_3$, so none of its actions can factor
through the abelian quotient.

For the exceptional packet the degree-two term vanishes on
$\Delta$, and odd terms vanish by G.2. Its first possible
nontrivial action is therefore in degree four on $N_5$;
G.7 proves that this action is nonzero. Again all higher
levels surject onto $N_5$. To specify the field without an
unnamed polylogarithmic class, choose compatible
$\zeta_{4\cdot5^n}$ with fourth-root component $i$, and put
\[
 \eta_n=\prod_{a\in(\mathbb Z/5^n)^\times}
          \sigma_a(1-\zeta_{4\cdot5^n})^a,
 \quad
 \sigma_a(i)=i,\quad
 \sigma_a(\zeta_{5^n})=\zeta_{5^n}^{a}.
\]
The conductor-four Soul\'e cocycle in
\citet[Definition 2.14 and Remark 2.15(2)]{IshiiGoto2026}
is characterized on the cyclotomic kernel by the Kummer
action on $\eta_n^{1/5^n}$. Consequently its restriction
to $\Delta$ cuts out the compositum
$K^{\rm ab}(\eta_n^{1/5^n}:n\ge1)$. The unit factors in
(G.8) leave its kernel unchanged. This also shows why an
outer elliptic Soul\'e invariant is the first obstruction
only after the central discriminant Kummer term has been
removed.
\end{proof}

\begin{statement}{G.18}{Known}{The class-field-theoretic CM symmetry quotient}
With the Artin reciprocity convention fixed, the quotient
$G_K^{\rm ab}$ is the finite-idelic class group
$\mathbb A_{K,f}^{\times}/K^{\times}$, and
$G_K^\circ/\Delta$ is the kernel of the induced cyclotomic
norm character to $\widehat{\mathbb Z}^{\times}$.
This is the symmetry group occurring in the CM system of
Connes, Marcolli, and Ramachandran. It describes the action
on $N_2$; it does not describe the genuine actions on
$N_r$ for $r\ge3$ at the specified tangent.
\end{statement}
\begin{proof}
For an imaginary quadratic field the connected archimedean
factor disappears from the Artin quotient, leaving the
displayed finite-idelic group. Norm compatibility of Artin
reciprocity identifies the cyclotomic character with the
norm followed by
$\mathbb A_{\mathbb Q,f}^{\times}/\mathbb Q_{>0}^{\times}
\simeq\widehat{\mathbb Z}^{\times}$, with inversion if the
opposite reciprocity convention is chosen. Taking kernels
proves the second assertion. The identification with the
CM dynamical symmetry is the assertion of
\citet[Theorem 5.1, source label \texttt{KMSCM}]{ConnesMarcolliRamachandran2005}
in the supplied source. G.17 proves the final limitation
by an explicit nonzero cocycle on the kernel of this
quotient, not by a difference of group names.
\end{proof}

\subsection{The remaining full profinite comparison}

\begin{statement}{G.19}{Proved}{Full profinite recognition when one prime detects the subgroup}
Let $\pi=\pi_1(U_{\overline K})$ be the full profinite geometric
fundamental group, and let $C\subset\pi$ be a nontrivial closed
procyclic subgroup with open arithmetic normalizer. If the image of
$C$ in $P_\ell$ is nontrivial for even one prime $\ell$, then $C$
is contained in a conjugate of actual full profinite cusp inertia.
\end{statement}
\begin{proof}
The $\ell$-component $C_\ell\simeq\mathbb Z_\ell$ maps injectively
onto its nonzero image in $P_\ell$: a nonzero homomorphism between
closed subgroups of $\mathbb Z_\ell$ is injective. G.12--G.13 show
that the normalizer's character on this image, hence on the actual
subgroup $C_\ell$ of $\pi$, is exactly $\chi_\ell$. Apply
\citet[Theorem 2.1.1]{Nakamura1994} with the full class of finite
groups to a generator of $C_\ell$. Its cyclotomic-normalizer
condition refers to the automorphism group of that generated
subgroup, here $\mathbb Z_\ell^\times$; it does not require the
subgroup to have a nonzero component at every prime. The theorem
puts $C_\ell$ inside one actual full inertia group $I$.

We justify explicitly why this controls all other components of
$C$. For any $1\ne c^a\in I$, including an exponent supported
only at $\ell$, its centralizer in $\pi$ is $I$. To see this,
choose an index-two geometric subgroup $H$ of the punctured-torus
group. It contains $c$, and $c$ is a member of a free profinite
basis of $H$, since the corresponding double cover has two
cusps. If $w$ centralizes $c^a$, then $w^2\in H$ also centralizes
it. \citet[Lemma 2.1.2]{Nakamura1994}, which explicitly allows
an arbitrary nonzero exponent in $\widehat{\mathbb Z}$, gives
$w^2\in I$. Since $c$ has zero image in $\pi^{\rm ab}$, it
follows that $2w^{\rm ab}=0$. That abelianization is the
torsion-free module $\widehat{\mathbb Z}^2$, so $w^{\rm ab}=0$
and $w\in H$. Apply the same lemma again to obtain $w\in I$.
This is the index-two argument of
\citet[\S1.3, Lemma 1]{Asada2001}, with the exponent generality
supplied by Nakamura's lemma. Finally $C$ is abelian and centralizes
a nonzero element of $C_\ell$, so $C\subset I$. No assertion
about its other power characters was used before this last step.
\end{proof}

\begin{statement}{G.20}{Open}{The precise remaining full profinite obligations}
Put $R=\bigcap_\ell\ker(\pi\to P_\ell)$. The full profinite
assertion in G(d) is reduced by G.19 to a procyclic subgroup
$1\ne C\subset R$ with open arithmetic normalizer. Its existence
or exclusion is undecided here. For G(b), G.10 proves that every
fixed element has inertia image in every $P_\ell$, but the
required full equality $\pi^{s_b(\Delta)}=I$ is not proved.
After removing the uniquely specified inertia element with those
primewise images, its exact remaining case is
\[
                         R^{s_b(\Delta)}=1.
\]
These are remaining obligations, not counterexamples or claimed
consequences of the pro-$\ell$ calculations.
\end{statement}
\begin{proof}[Attacks, reductions, and exact obstruction]
The primewise images of a fixed $u$ are $c^{a_\ell}$ by G.10.
The exponent $(a_\ell)_\ell\in\widehat{\mathbb Z}$ gives a
single actual inertia element. Multiplying its inverse by $u$
produces a fixed element of $R$, proving the stated reduction.
The first attack was therefore to use all nilpotent and pro-$\ell$
quotients together. Those quotients do not separate the full free
profinite group, so this gives no proof that the remaining element
is trivial.

The second attack uses finite covers. Any nonidentity $g\in\pi$
has nonzero abelian image in some open subgroup: take a finite
quotient in which its image $\bar g$ is nontrivial and pull back
$\langle\bar g\rangle$. The quotient onto that cyclic group
factors through the open subgroup's abelianization and detects $g$.
After finite extension the corresponding geometric cover descends,
and an open arithmetic subgroup stabilizes it. For a
$\Delta$-fixed element this gives a nonzero class fixed by
$\Delta\cap G_{K'}$ in the cover's homology. The missing assertion
is that this class must lie in the boundary part. G.22 narrows
the alternative: its proper homology image lies in the Prym
kernel of $J_{\overline Y}\to E$, so the inherited elliptic
quotient cannot carry this residual vector, and a genus-one
compactification is excluded. A higher-genus Prym may still
have $\Delta$-fixed CM components. The argument cannot discard
them merely by invoking Frobenius purity, because $\Delta$ is
of infinite index and fixes the relevant CM representations.
A sufficient additional
step would be a cofinal construction of detecting covers for which
the detected vector avoids every such fixed CM summand. No such
construction has been established in this return.

The third attack checks nonlinear inertia-recognition theorems.
Nakamura's Theorem 2.1.1 does solve G.19 once a cyclotomic
prime component is obtained. For a subgroup wholly in $R$, the
lower-central eigencharacter argument has no nonzero vector on
the original curve. Passing to an open cover restores a vector,
but its character can have Jacobian weight one; that is precisely
the unresolved alternative. Applying the theorem over $K^{\rm ab}$
instead is not justified. The extension in
\citet[\S5, Theorem 20 and Corollary 22]{Stix2008Cuspidal}
requires the base field to distinguish tori from abelian varieties.
Over $K^{\rm ab}$ both $\mathbb Z_\ell(1)$ and $T_\ell E$
are trivial representations, so this hypothesis fails explicitly.
The cyclotomic-kernel results CK A6--A7 consequently do not cover
this subgroup. The peripheral-centralizer theorem alone also does
not place a fixed element, or an undetected normalizer, in inertia.

The fourth attack tries to send a detected Prym vector to a
punctured CM elliptic curve. G.22 proves the exact conditional
transport statement. A nonconstant $f:\overline Y\to E'$ must
have an entire fibre supported in $S=\overline Y\setminus Y$;
this forces $\deg f\le |S|+2g(\overline Y)-2=\deg(Y/U)$.
A detecting Jacobian quotient does not ensure this condition:
$[3]:E\to E$ detects proper homology of the degree-four cover
$E\setminus E[2]\to U$, but its restriction still surjects
onto all of $E$. Further covers do not repair that fixed map.
Deleting extra points instead requires an equivariant lift of
the subgroup, and G.22 gives an explicit CM eigenline for which
an open-normalizer lift fails. Even existence of an elliptic
quotient, and the recognition theorem for its exact pointed
target, are additional requirements for a general CM Prym
factor. Thus the missing step is now a detecting map with this
fibre condition to a target whose normalizer criterion is proved,
or a different argument excluding the residual Prym alternative.
The fixed-element version must in addition transport the actual
lifted cusp-packet section. These requirements have not been
established for arbitrary detecting covers.
\end{proof}

\begin{statement}{G.22}{Proved}{Prym detection and the condition for punctured elliptic transport}
Let $k/K$ be finite and let $p:Y\to U_k$ be a geometrically
connected finite \emph{\'etale} cover of degree $d$. Write
$\overline p:\overline Y\to E_k$ for its extension to smooth
proper curves, $S=\overline Y\setminus Y$, and
$g_Y=g(\overline Y)$; cardinalities of boundary sets are geometric.
The following assertions delimit the finite-cover attack on G.20.
\begin{enumerate}
\item If $u\in\pi_1(Y_{\overline k})$ has image in $R$, its
proper homology class $\overline v_\ell$ lies in
\[
 \ker\bigl(V_\ell J_{\overline Y}
              \xrightarrow{\overline p_*}V_\ell E\bigr)
       =V_\ell\bigl(\ker(\overline p_*:J_{\overline Y}\to E)^0\bigr).
\]
In particular $\overline v_\ell=0$ if $g_Y=1$. A residual
weight-one detector must therefore occur in the Prym part of a
cover of genus at least two.
\item For a nonconstant $f:\overline Y\to E'$ and a geometric
point $q\in E'$, there is a restriction
$Y\to E'\setminus\{q\}$ exactly when
\[
                    \operatorname{Supp}f^{-1}(q)\subset S.
\]
This condition implies
\[
                 \deg f\le |S|+2g_Y-2=d.
\]
Equality holds exactly when $S=\operatorname{Supp}f^{-1}(q)$
and $f$ is unramified away from that fibre; then the restricted
map is finite \emph{\'etale}.
\item Suppose $f$ has this fibre condition and is defined, with
$q$, over a finite extension of $k$. Suppose also that the
pro-$\ell$ arithmetic normalizer criterion of G.13 holds for
the \emph{particular pointed curve} $(E',q)$. No procyclic
subgroup of $\pi_1(Y_{\overline k})$ with open arithmetic
normalizer can have proper homology image carried nontrivially
by $f_*$ to $V_\ell E'$. The target hypothesis is established
here when $(E',q)$ becomes isomorphic to $(E,O)$ after a finite
extension. It is not asserted for arbitrary isogenous curves,
CM fields, or higher-dimensional CM factors.
\item A detecting elliptic Jacobian quotient need not satisfy
the fibre condition. Explicitly, $[2]:E\setminus E[2]\to U$
is a degree-four finite \emph{\'etale} cover, and
$[3]:E\to E$ is injective on every Tate module, but
\[
                      [3](E\setminus E[2])=E.
\]
It therefore gives no map to any once-punctured $E$.
\item Adding the missing punctures does not in general permit
an arithmetic-normalizer-preserving lift of a procyclic subgroup.
A primitive CM eigenline in $T_5E$ has open arithmetic normalizer,
but no procyclic subgroup of $P_5$ with open arithmetic normalizer
maps nontrivially into that line under the filling homomorphism
$P_5\twoheadrightarrow T_5E$.
\end{enumerate}
\end{statement}
\begin{proof}
The map $\pi_1(Y_{\overline k})\to\pi$ is the geometric
inclusion determined by a basepoint path. The induced homomorphism
to the abelian pro-$\ell$ quotient $T_\ell E$ factors through
$T_\ell J_{\overline Y}$ and $\overline p_*$. It vanishes on
an element mapping into $R$, because this quotient factors
through $P_\ell$. Tensoring with $\mathbb Q_\ell$ proves the
kernel assertion. Factoring
the surjective homomorphism $J_{\overline Y}\to E$ by the
identity component of its kernel leaves an isogeny, which is
injective on Tate modules; this gives the displayed equality.
If $g_Y=1$, Riemann--Hurwitz makes $\overline p$ unramified.
After a finite extension and choices of origins it is an isogeny
of elliptic curves. Its Tate-module map is injective: a compatible
torsion sequence in its finite kernel must be zero, since a
fixed power of $\ell$ kills that kernel's $\ell$-part. This
proves the genus-one assertion. In particular the original
elliptic quotient of $J_{\overline Y}$ cannot itself carry the
residual vector; any relevant CM component must occur in the
indicated kernel.

The fibre condition is just the condition that the restriction
avoid $q$. All ramification of $\overline p$ lies above $O$,
and that fibre is supported on $S$. In characteristic zero,
Riemann--Hurwitz gives
\[
 2g_Y-2=\sum_{P\in S}(e_P(\overline p)-1)=d-|S|.
\]
For $f$ it gives $\sum_P(e_P(f)-1)=2g_Y-2$; see
\citet[\S53.12, Lemmas 53.12.2 and 53.12.4]{GStacksRiemannHurwitz}.
Under the fibre condition,
\[
 \deg f
   =|f^{-1}(q)|+\sum_{P\in f^{-1}(q)}(e_P(f)-1)
   \le |S|+2g_Y-2=d.
\]
Here the first cardinality counts distinct geometric points.
Equality requires equality in both inequalities, precisely as
stated; the resulting restriction is finite and unramified
between smooth curves, hence finite \emph{\'etale}.

For the third assertion, make the finite base extension on which
all maps are defined. An open arithmetic image stays open under
intersection with the corresponding open arithmetic fundamental
group. Choose compatible geometric basepoints and a path. The
morphism $Y\to E'\setminus\{q\}$ induces a homomorphism of
arithmetic fundamental groups over $G_k$, and hence one to the
geometrically pro-$\ell$ target. It carries a normalizer into
the normalizer of the image subgroup and preserves its open
arithmetic image. The nonzero proper homology image ensures
that the image subgroup is a nontrivial closed procyclic
subgroup of the target geometric group. The assumed normalizer
criterion puts it in cusp inertia. But cusp inertia has zero
image in $T_\ell E'$, contradicting the detecting hypothesis.
When the pointed target becomes $(E,O)$, this is exactly G.13:
the finite extension changes the arithmetic group to an open
subgroup, and changes of basepoint only conjugate the relevant
subgroups. No claim about a different CM curve follows from
this argument merely because its Jacobian has an abelian
Galois representation.

For application to an original normalizer, one can first choose
an open geometric cover containing its procyclic subgroup.
Intersecting its normalizer with an arithmetic fundamental
group of that descended cover leaves an open arithmetic image,
so the preceding functorial argument applies. The missing step
is a suitable $f$ satisfying both detection and the fibre
condition. A CM character on a Tate module alone supplies neither
a rational elliptic quotient nor this condition: a CM factor
can have dimension greater than one, and its individual
$\mathbb Q_\ell$ eigenlines need not come from elliptic
quotients.

For the explicit example, $E[2]=\{O,(0,0),(1,0),(-1,0)\}$,
and $[2]$ is unramified in characteristic zero. Every fibre of
$[3]$ has nine geometric points. Since $[3]$ acts as the
identity on $E[2]$, a fibre over $q\in E[2]$ contains exactly
one point of $E[2]$ and eight points of its complement. A fibre
over $q\notin E[2]$ is wholly in the complement. This proves
the displayed surjectivity. The proper homology map is $3I$,
which is injective even over $\mathbb Z_3$.
The exact independent certificate in
\texttt{code/g\_cover\_transport.py} derives
\[
 \psi_3(x)=3x^4-6x^2-1,
 \qquad x(2P)-x(P)=\frac{-\psi_3(x(P))}{4(x(P)^3-x(P))}.
\]
It verifies that $\psi_3$ is squarefree and relatively prime
to $x^3-x$, and that its values at $0,1,-1$ are $-1,-4,-4$.
Thus its four roots give eight nonzero points of order three,
all outside $E[2]$, in agreement with the fibre argument.
This example refutes only automatic puncturing of a detecting
quotient; it is not a residual element or a counterexample to
G.20. More generally, postcomposing any detecting $f$ by $[m]$
preserves nonzero Tate-module detection but multiplies its
degree by $m^2$. For large $m$ the degree bound excludes every
puncture. Passing to a further surjective cover of $Y$ cannot
repair a fixed map whose image is all of $E'$.

Finally, CM over $K$ and the splitting of $5$ give a primitive
$\mathbb Z_5$ eigenline $L\subset T_5E$ stable under $G_K$.
The zero section splits the arithmetic fundamental group of the
proper elliptic curve, so the normalizer of $L$ has full
arithmetic image. If a closed procyclic subgroup upstairs in
$P_5$ mapped nontrivially to $L$ and retained an open arithmetic
normalizer, G.13 would make it cuspidal. Filling the puncture
kills its image, a contradiction. This failure persists after
any finite base extension. Likewise G.10 shows that a nonzero
$\Delta$-fixed vector of $T_5E$ has no lift fixed by any of
the specified cusp-packet actions on $P_5$. Thus deleting more
points cannot be used as an automatic equivariant lifting step.
The conditional transport excludes a residual normalizer
\emph{when} the required map exists; it does not establish that
existence for the covers in G.20.
\end{proof}

